% -*- mode:LaTex; mode:visual-line; mode:flyspell; fill-column:75-*-

\chapter{Introduction} \label{chapIntroduction}

The U.S. apple industry is a cornerstone of American agriculture, contributing
more than \$23 billion to the national economy and supporting over 150,000
jobs across a supply chain that spans orchards, packing houses, cold storage,
and distribution~\cite{gerlach2025usapple}.
Yet despite record production levels, many growers are operating at a loss,
and rising labor costs are making the intensive manual scouting that disease
management depends on increasingly difficult to
sustain~\cite{gerlach2025usapple}.
In this environment, disease management is not merely an agronomic
concern: it is an economic one, since a single missed outbreak can consume
a season's already thin margin.

Fire blight, caused by the bacterium \textit{Erwinia amylovora}, is among the
most destructive threats facing apple growers today.
Annual losses and management costs in the United States exceed \$100 million,
and individual outbreaks can devastate entire orchards in a single
season~\cite{Norelli2003FireBlightManagementTwentyFirst,
acimovic2026fireblight}.
The bacterium overwinters in bark lesions called cankers and spreads in spring
to infect blossoms and young shoots~\cite{van1979FireBlight}; infected shoots
develop a characteristic ``shepherd's crook'' deformation, while infected
woody tissue forms new cankers (Figure~\ref{figFireBlightSymptoms}) and, in
severe cases, girdles and kills the rootstock or trunk outright.
Effective management depends on early detection of these cankers and
shepherd's crooks during the dormant season, yet current practice relies
almost entirely on manual scouting, a labor-intensive, subjective process
that does not scale to the multi-acre plantings of modern commercial
operations and is increasingly constrained by a shrinking agricultural
workforce~\cite{gerlach2025usapple}.

Robotic monitoring systems offer a promising path forward.
Advances in vision-based detection, manipulator-based active sensing, and
probabilistic three-dimensional mapping have made it possible to deploy
robotic systems for sustained, repeatable inspection of individual
plants~\cite{droukas2023survey}.
The orchard canopy, however, remains a challenging environment for autonomous
perception: canopy variation, severe occlusion, thin branch geometry, and
large swings in natural lighting all degrade the reliability of fixed-camera,
fixed-viewpoint sensing.
The visual cues that distinguish infected tissue from healthy tissue are
subtle, and resolving them reliably requires close-range imaging at fine
spatial resolution, conditions that fixed, wide-angle views cannot
consistently provide.
Moreover, we aim to provide technology for growers capable of branch-level
disease localization, which will enable growers to assess the severity of
disease and make informed decisions about pruning and treatment.

% Despite growing research interest in agricultural robotics, and although
% close-range vision has recently been shown to detect dormant cankers on
% pear~\cite{Linker2024MappingFireBlightCankersFRCNN}, the specific problem of
% autonomously detecting and spatially localizing fire blight cankers and
% shepherd's crooks throughout the canopy of dormant apple trees has received
% little attention.

% TODO: Add figs/robot_symptoms_collage_2.png, then uncomment.
% \begin{figure}[t]
%     \centering
%     \includegraphics[width=0.8\linewidth]{figs/robot_symptoms_collage_2.png}
%     \caption{Autonomous disease inspection system (left) and representative
% fire blight symptoms. Shoot blight instances exhibit a characteristic
% shepherd's crook (top right), while canker regions appear as sunken or
% darkened lesions on woody tissue (bottom right).}
%     \label{figSystemAndSymptoms}
% \end{figure}

These limitations motivate an active perception approach in which a robotic
manipulator positions a camera to acquire discriminative, task-relevant views.
This approach takes advantage of specialized sensors, namely flash RGB and
near-infrared cameras, that are highly effective for disease detection because
they can resolve the fine bark texture characteristic of cankers, but that
correspondingly have narrow fields of view.
Realizing this approach end-to-end requires three components that, to our
knowledge, have not previously been brought together for fire blight: a
dataset and detector capable of recognizing disease symptoms across sensing
modalities, a three-dimensional representation that can accumulate and
reconcile noisy per-view detections into a single map of the canopy, and a
planner that uses that representation to decide where to look next. This
thesis aims to address these gaps directly.

% TODO: Add figs/system_diagram.png, then uncomment.
% Diagram should show three linked parts: (1) hardware -- Amiga mobile base
% with xArm6 manipulator positioned in front of a dormant apple tree, camera
% payload (flash stereo RGB + NIR) on the end effector; (2) the closed
% perception--action loop -- camera capture -> YOLO-seg 2D detector ->
% fusion into the confidence-aware semantic occupancy map -> NBV planner
% selects next viewpoint -> arm moves -> back to capture; (3) output --
% inset of the resulting semantic map (disease-confidence-colored occupancy
% grid) showing the spatial confidence map delivered to the grower.
% \begin{figure}[t]
%     \centering
%     \includegraphics[width=\linewidth]{figs/system_diagram.png}
%     \caption{Overview of the proposed active perception system. The Amiga
% mobile base positions the xArm6 manipulator in front of a dormant apple
% tree; the arm iteratively selects and executes inspection viewpoints,
% capturing flash stereo RGB and NIR multispectral imagery. Per-view
% detections from the YOLO-seg detector are fused into a confidence-aware
% semantic occupancy map, which the NBV planner queries to select the next
% viewpoint, closing the perception--action loop.}
%     \label{figSystemDiagram}
% \end{figure}


\section{Contributions} \label{secContributions}

The primary contributions of this thesis are organized around three
components: a multi-modal dataset and YOLO-based fire blight detector, a
semantic mapping and planning framework, and the field-validated robotic
system that integrates them.

\subsection{Multi-Modal Dataset and Fire Blight Detection}

We collect a dataset of flash stereo RGB, 25-band NIR multispectral, and ZED2i
RGB-D imagery of dormant apple trees exhibiting a range of fire blight
symptoms, gathered across multiple collection sessions in orchards in Virginia
and the greater Pittsburgh region.
Manual expert annotation is accelerated with a SAM-assisted labeling workflow,
and a cross-modal label transfer procedure propagates RGB annotations onto the
corresponding NIR frames; to our knowledge, this is the first dataset pairing
dense NIR spectral bands with RGB imagery for fire blight detection on dormant
woody tissue.
Building on this dataset, we train YOLO detection models and compare
performance between the RGB and NIR modalities.
The dataset, annotation methodology, and detection results are presented in
Chapter~\ref{chapDataset}.

\subsection{Semantic Mapping and Active Viewpoint Selection}

We introduce a confidence-aware semantic occupancy representation that jointly
maintains occupancy probability, semantic disease labels, and prediction
confidence within a unified volumetric map, fusing per-view detections into a
single, persistent representation of disease confidence across the canopy.
Built on this representation, we develop an active viewpoint selection
framework centered on a semantic-aware next-best-view (NBV) planner that
prioritizes low-confidence symptomatic regions for re-observation, evaluated
alongside geometric coverage-based and volumetric NBV baselines to enable
refined observation of complex canopy structures.
The mapping representation, the planner, and their integration are presented
in Chapter~\ref{chapNBV}.

\subsection{A Field-Validated Robotic System}

We combine these components into Erwin, a robotic active perception system
for targeted fire blight detection in dormant apple trees, built around an
Amiga mobile base and a UFACTORY xArm6 manipulator.
The mobile base remains fixed in front of a tree during inspection while the
arm autonomously selects and executes inspection viewpoints within a closed
perception--action loop, combining flash stereo RGB and 25-band near-infrared
multispectral imaging with the confidence-aware semantic occupancy map and NBV
planner described above.
Erwin is deployed and validated on real apple trees at the Penn State Fruit
Research and Extension Center in Adams County, Pennsylvania, with mapping
performance assessed against splat-based 3D reconstructions.
Validation results are presented in Chapter~\ref{chapErwin}.
The work presented here focuses on per-tree detection, mapping, and active
viewpoint selection; extending the system to traverse and map an entire
orchard autonomously, row by row, is an important direction for future work
(Chapter~\ref{chapConclusions}).
Erwin was developed as part of Carnegie Mellon University's Fire Blighters
team entry in the 2026 Farm Robotics Challenge, where the team received the
Amiga Innovation Award.

Together, these contributions form an end-to-end pipeline: the dataset and 2D
detector supply the per-view symptom evidence consumed by the confidence-aware
semantic occupancy map and NBV planner, which are in turn integrated onto
Erwin and validated under field conditions.

\section{Thesis Organization} \label{secOutline}

The remainder of this thesis is organized as follows.
Chapter~\ref{chapBackground} provides background on agricultural robotics and
active perception, fire blight biology and current management practice, and
prior work on automated fire blight and orchard disease diagnosis.
Chapter~\ref{chapDataset} describes the sensor rig, the data collection
procedure, the multi-modal orchard dataset, and the 2D detection pipeline
trained on it.
Chapter~\ref{chapNBV} presents the confidence-aware semantic occupancy
representation and the proposed NBV planners.
Chapter~\ref{chapErwin} describes the full integrated robotic system, Erwin,
and reports field validation results.
Chapter~\ref{chapConclusions} concludes the thesis with a discussion of
limitations and future directions.
